\documentclass[11pt]{article}
\usepackage[margin=1in]{geometry}
\usepackage{amsmath,amssymb,amsthm}
\usepackage{graphicx}
\usepackage{booktabs}
\usepackage{hyperref}
\usepackage{xcolor}
\usepackage{natbib}
\usepackage{subcaption}
\usepackage{multirow}

\newcommand{\citeneeded}[1]{\textcolor{red}{[CITE: #1]}}
\newcommand{\todo}[1]{\textcolor{orange}{[TODO: #1]}}

\title{Dynamical composition creates and destroys epistasis:\\
order-by-order quantification in six Boolean gene regulatory networks}

\author{
  Elliot Tower \\
  \texttt{elliot@elliottower.ai}
}

\begin{document}
\maketitle

\begin{abstract}
Nonlinear dynamics generate phenotypic epistasis from molecular
architectures that may contain no epistatic interactions themselves.
Prior work has established this gap qualitatively; we quantify it
order-by-order. For six published Boolean gene regulatory networks
(10--15 genes), we compute the complete Walsh--Hadamard interaction
spectrum at two levels: (i)~the local update rules, where each gene's
truth table yields exact Fourier coefficients, and (ii)~the
attractor-level knockout landscape, obtained from exhaustive $2^n$
coalition sweeps under synchronous dynamics. Comparing the two spectra
reveals a signed \emph{composition gap} whose magnitude spans two
orders: from effectively zero in the fission yeast cell cycle to
$+34.3$ percentage points of order-3+ energy created in
\emph{Arabidopsis}. Pairwise Spearman correlation between local and
global interaction magnitudes is universally positive
($\rho = 0.12$--$0.41$) but statistically significant in only three of
six networks, indicating that molecular wiring is a noisy predictor of
knockout epistasis. A pre-registered cross-network hypothesis---that
positive feedback loop ratio predicts the gap direction---fails
decisively. An exploratory association between limit-cycle prevalence
and gap direction is noted ($p \approx 0.10$, post-hoc) as a
hypothesis for future testing. The complete Walsh spectra, coalition
tables, and pre-registration artifacts are released for reuse.
\end{abstract}


%% ====================================================================
\section{Introduction}
\label{sec:intro}
%% ====================================================================

Epistasis---the dependence of phenotype on gene combinations rather
than individual effects---is among the oldest measurement problems in
quantitative genetics \citeneeded{Fisher 1918, Bateson 1909}. A
persistent difficulty is the gap between molecular and statistical
epistasis: the interactions encoded in regulatory logic need not
predict the interactions observed in knockout experiments
\citeneeded{Phillips 2008 Nat Rev Genet}. Gjuvsland et al.\ showed
that nonlinear dynamics generically produce statistical epistasis from
gene regulatory architectures, regardless of whether the underlying
molecular rules contain epistatic terms \citeneeded{Gjuvsland 2007
Genetics}. Subsequent work confirmed this in synthetic three-gene
circuits, where pairwise and triplet epistasis depended on environment
and switched sign across conditions \citeneeded{Baier 2023 Sci Adv}.

These results establish that a gap exists. They do not quantify it.
The gap has no known magnitude, no known sign structure, and no
established dependence on network architecture. One reason is
computational: exhaustive knockout measurement requires $2^n$ subset
evaluations, which is infeasible at genomic scale. In Boolean gene
regulatory networks of modest size ($n \leq 15$), however, exhaustive
evaluation is tractable.

We exploit this tractability. For each network, we compute two
complete interaction spectra using the Walsh--Hadamard transform: one
from the local update rules (each gene's truth table, at most $2^6 =
64$ entries) and one from the attractor-level value function (all
$2^n$ knockout coalitions, up to $2^{15} = 32{,}768$ evaluations).
The two spectra are compared order-by-order, yielding a signed
composition gap that measures how much higher-order structure dynamics
create or destroy.

The central finding is that the composition gap is large, signed, and
architecture-dependent. Dynamics create higher-order epistasis in some
networks (adding interactions absent from any local rule) and destroy
it in others (compressing higher-order structure into lower orders),
with one network producing a near-zero gap. The magnitudes span two
orders, from sub-percentage-point to 34 percentage points.
A pre-registered hypothesis based on positive feedback loop
counting fails, and the molecular wiring diagram is a statistically
significant predictor of knockout epistasis in only half the networks
studied.


%% ====================================================================
\section{Methods}
\label{sec:methods}
%% ====================================================================

\subsection{Boolean network models}
\label{sec:models}

Six published Boolean gene regulatory networks were selected from the
Cell Collective repository \citeneeded{Helikar 2012} and primary
literature (Table~\ref{tab:models}). Selection criteria: (i)~$n \leq
15$ nodes to permit exhaustive coalition sweeps, (ii)~synchronous
update semantics used in the original publication, and
(iii)~biologically defined output node (cell-cycle commitment marker,
apoptosis effector, or DNA repair checkpoint). Three networks model
cell-cycle control across species (mammalian, \emph{Drosophila},
\emph{Arabidopsis}), one models apoptosis signaling, one models
fission yeast cell-cycle, and one models the Fanconi anemia DNA
repair pathway.

\begin{table}[t]
\centering
\caption{Boolean gene regulatory networks used in this study. $n$:
number of genes (players in the coalition game). Output: the gene
whose attractor state defines the value function. Feedback loops
counted up to length~4.}
\label{tab:models}
\begin{tabular}{llrccrr}
\toprule
ID & Network & $n$ & Output & Pos.\ loops & Neg.\ loops & $2^n$ \\
\midrule
G1 & Faure mammalian cell cycle & 10 & CycB & 13 & 12 & 1{,}024 \\
G2 & Tournier apoptosis & 12 & C3a & 3 & 3 & 4{,}096 \\
G4 & Davidich fission yeast & 10 & Cdc2\_Cdc13 & 13 & 3 & 1{,}024 \\
G5 & Albert \emph{Drosophila} cell cycle & 14 & CycB & 22 & 14 & 16{,}384 \\
G9 & Rodriguez Fanconi anemia & 15 & CHKREC & 36 & 45 & 32{,}768 \\
G10 & Ortiz-Guti\'errez \emph{Arabidopsis} cell cycle & 14 & CYCB1\_1 & 39 & 42 & 16{,}384 \\
\bottomrule
\end{tabular}
\end{table}


\subsection{Local rule Fourier extraction}
\label{sec:local-fourier}

Each gene $k$ has a Boolean update rule $f_k\colon \{0,1\}^{r_k}
\to \{0,1\}$, where $r_k \leq 6$ is the number of regulators. We
compute the Walsh--Hadamard transform of the truth table:
%
\begin{equation}
\label{eq:local-wht}
w_T^{(k)} = \frac{1}{2^{r_k}} \sum_{x \in \{0,1\}^{r_k}} f_k(x)\,(-1)^{|x \cap T|},
\quad T \subseteq \{1, \ldots, r_k\}.
\end{equation}
%
A nonzero coefficient $w_T^{(k)}$ at $|T| = m$ indicates a genuine
$m$-way interaction among the regulators of gene $k$,
representation-independent: the transform is invariant to how the
rule is written (disjunctive normal form, algebraic normal form, etc.).
For each unordered pair $\{i,j\}$ of network genes, the
\emph{local pairwise magnitude} is:
%
\begin{equation}
\label{eq:local-pair}
\ell_{ij} = \sum_{k:\, i,j \in \mathrm{reg}(k)} |w_{\{i,j\}}^{(k)}|,
\end{equation}
%
summing over all genes $k$ whose update rule depends on both $i$ and
$j$.

The \emph{local energy spectrum} reports the fraction of total
squared-coefficient energy at each interaction order $m$:
%
\begin{equation}
\label{eq:local-spectrum}
E_m^{\mathrm{local}} = \frac{\sum_k \sum_{|T|=m} \bigl(w_T^{(k)}\bigr)^2}
{\sum_k \sum_T \bigl(w_T^{(k)}\bigr)^2}.
\end{equation}


\subsection{Coalition sweep and attractor-level value function}
\label{sec:coalition-sweep}

For each of $2^n$ coalitions $S \subseteq \{1,\ldots,n\}$, genes not
in $S$ are clamped to~0 (knockout). The remaining genes evolve under
synchronous Boolean dynamics from $N = 512$ uniformly random initial
states until convergence to a fixed point or limit cycle (detected by
Floyd's algorithm, maximum 200 steps). The value function is:
%
\begin{equation}
\label{eq:value}
v(S) = \frac{1}{N} \sum_{i=1}^{N}
\overline{x_{\mathrm{out}}^{(i)}(S)},
\end{equation}
%
where $\overline{x_{\mathrm{out}}^{(i)}(S)}$ is the time-averaged
output-node state for initial condition $i$ under coalition $S$ (for
fixed points, the attractor value; for limit cycles, the mean over
one period).

The global Walsh--Hadamard transform of $v$ gives the
\emph{global interaction spectrum}:
%
\begin{equation}
\label{eq:global-wht}
w_T^{\mathrm{global}} = \frac{1}{2^n} \sum_{S \subseteq \{1,\ldots,n\}}
v(S)\,(-1)^{|S \cap T|}.
\end{equation}
%
The global pairwise magnitude for pair $\{i,j\}$ is
$g_{ij} = |w_{\{i,j\}}^{\mathrm{global}}|$, and the
\emph{global energy spectrum} is defined analogously to
Eq.~\ref{eq:local-spectrum}.


\subsection{Composition gap metrics}
\label{sec:metrics}

\textbf{Pairwise Spearman correlation.}
For each network, we compute
$\rho = \mathrm{Spearman}(\ell_{ij}, g_{ij})$ over all
$\binom{n}{2}$ pairs, with bootstrap 95\% confidence intervals
(10{,}000 resamples of initial states). A positive $\rho$ means local
rule structure is informative about global knockout epistasis; $\rho
\approx 0$ means the molecular wiring diagram provides no useful
prior.

\textbf{Order-3+ energy gap.}
The \emph{composition gap} at higher orders is:
%
\begin{equation}
\Delta_{3+} = \sum_{m \geq 3} E_m^{\mathrm{global}}
- \sum_{m \geq 3} E_m^{\mathrm{local}}.
\end{equation}
%
Positive $\Delta_{3+}$ indicates \emph{creation}: dynamics generate
higher-order interactions absent from any single local rule. Negative
$\Delta_{3+}$ indicates \emph{destruction}: dynamics compress
higher-order structure into lower orders.

\textbf{Quality gate.}
The global spectrum is analyzed only if (i)~the value function has
$\geq 10$ unique values (rounded to $10^{-6}$), (ii)~absolute
order-3+ energy exceeds $10^{-4}$, and (iii)~no single value accounts
for $> 90\%$ of coalitions. All six networks pass this gate.


\subsection{Blind prediction protocol}
\label{sec:blind}

To test whether structural features predict the composition gap
\emph{a priori}, a structurally-blinded analysis was conducted before
any coalition sweep ran. For each network, a separate analysis
examined the Boolean rules, feedback loop topology (up to length~4),
AND-gate arity, pathway depth, functional redundancy, and output-node
rule structure. From this analysis alone, the following predictions
were registered (SHA-frozen in \texttt{blind\_predictions.json}):
%
\begin{enumerate}
\item Spearman $\rho$ range (e.g., $[0.2, 0.5]$) and direction
\item Creation vs.\ destruction of higher-order interactions
\item Global energy spectrum ranges
\end{enumerate}
%
A cross-network hypothesis was also registered: the positive feedback
loop ratio
$\mathrm{fb\_ratio} = n_{\mathrm{pos}} / (n_{\mathrm{pos}} +
n_{\mathrm{neg}})$ predicts the gap direction ($> 0.5$ implies
creation).


%% ====================================================================
\section{Results}
\label{sec:results}
%% ====================================================================

\subsection{The composition gap is large, signed, and universal}

Table~\ref{tab:results} summarizes all six networks. Two findings
hold across the full panel.

\begin{table}[t]
\centering
\caption{Composition gap results for six Boolean GRNs. $\rho$:
Spearman correlation between local and global pairwise magnitudes
(bootstrap 95\% CI). $E_{3+}$: fractional energy at interaction
order $\geq 3$. $\Delta_{3+}$: global minus local $E_{3+}$.
Cycling: fraction of (coalition $\times$ initial state) pairs
reaching limit cycles.}
\label{tab:results}
\begin{tabular}{lrcccccc}
\toprule
Network & $n$ & $\rho$ [95\% CI] & $p$ & Global $E_{3+}$ & Local $E_{3+}$ & $\Delta_{3+}$ & Cycling \\
\midrule
G10 \emph{Arabidopsis} & 14 & 0.205 [$-$0.01, 0.41] & 0.052 & 38.3\% & 4.0\% & $+$34.3 pp & 16.4\% \\
G2 Tournier apoptosis & 12 & 0.247 [0.03, 0.45] & 0.046 & 12.6\% & 4.2\% & $+$8.5 pp & 7.4\% \\
G5 \emph{Drosophila} & 14 & 0.120 [$-$0.09, 0.32] & 0.257 & 10.4\% & 4.3\% & $+$6.2 pp & 14.0\% \\
G9 Fanconi anemia & 15 & 0.305 [0.12, 0.48] & 0.002 & 1.0\% & 6.5\% & $-$5.5 pp & 47.5\% \\
G1 Faure cell cycle & 10 & 0.409 [0.13, 0.64] & 0.005 & 0.8\% & 3.8\% & $-$3.0 pp & 19.6\% \\
G4 Davidich yeast$^\dagger$ & 10 & 0.154 [$-$0.15, 0.42] & 0.312 & 2.2\% & 2.6\% & $-$0.4 pp & 17.7\% \\
\bottomrule
\multicolumn{8}{l}{\footnotesize $^\dagger$G4 $\Delta_{3+}$ is 8$\times$ smaller than the next-smallest; effectively neutral.}
\end{tabular}
\end{table}

\textbf{Universally positive correlation.}
All six networks show positive Spearman $\rho$ between local and
global pairwise magnitudes ($\rho = 0.12$--$0.41$). Local rule
structure is informative about global knockout epistasis in every
network studied. The correlation is statistically significant in three
networks (G1, G2, G9; $p < 0.05$) and not significant in three (G4,
G5, G10). The confidence intervals overlap substantially across all
six, consistent with a shared underlying $\rho$ near $0.2$--$0.3$.

\textbf{Signed composition gap.}
Three networks show clear creation of higher-order interactions
($\Delta_{3+} = +6$ to $+34$~pp), two show clear destruction
($\Delta_{3+} = -3$ and $-5.5$~pp), and one (G4) is effectively
neutral ($\Delta_{3+} = -0.4$~pp, $8\times$ smaller than the
next-smallest effect). The magnitudes span two orders of magnitude,
establishing the composition gap as a large effect, comparable to
the local higher-order energy itself. G10 (\emph{Arabidopsis})
dominates the creation side: its $+34.3$~pp gap is $4\times$ the
runner-up (G2, $+8.5$~pp) and larger than the remaining five networks'
effects combined.


\subsection{The \emph{Arabidopsis} outlier}
\label{sec:arabidopsis}

G10 (\emph{Arabidopsis} cell cycle) produces the largest composition
gap by an order of magnitude. Only 9.6\% of global value function
energy remains at order~0, compared to 55.4\% in the local rules.
Dynamics redistribute constant and linear energy almost entirely into
higher-order interactions: 38.3\% of global energy is at order~$3+$,
a $9.6\times$ amplification from the local 4.0\%.

Several structural features distinguish G10. The network has the
highest edge density in the panel (63 edges for 14 nodes, 4.5
edges/node). The output gene CYCB1\_1 depends on 8 regulators through
deeply nested AND/OR logic, and E2Fe---which feeds into the output
pathway---has 6 regulators with maximal Boolean nesting. These
features create combinatorial pathway compounding: perturbations
propagate through dense interconnections with many opportunities for
interaction at each step.

Despite the massive higher-order creation, pairwise Spearman
correlation is only marginally non-significant ($\rho = 0.205$,
$p = 0.052$). Triple-level correlation is highly significant
($\rho = 0.163$, $p = 0.002$), suggesting that the local--global
relationship operates more strongly at higher interaction orders in
dense networks.

Because G10 is $4\times$ the next-largest creation effect, any
aggregate statistic for ``creation networks'' is effectively a
statement about \emph{Arabidopsis}. This network should be treated
as a case study in extreme composition gap amplification, driven by
exceptional edge density and output-rule complexity, rather than as
a representative data point.


\subsection{Creation and destruction mechanisms}
\label{sec:mechanisms}

The three creation networks share a structural motif: complex output
gating through multiple independent regulatory paths.
%
\begin{itemize}
\item G10: CYCB1\_1 gated by 8 regulators with nested AND/OR.
\item G2: Output C3a requires TNF signaling, insufficient survival
  signaling, and C8a activation---a three-way conjunction.
\item G5: Output CycB gated by Fzy, Fzr, and Wee1/Stg, fed by
  three independent input nodes (Ago, CycD, Notch).
\end{itemize}
%
The two clear destruction networks share simple output rules.
%
\begin{itemize}
\item G1: Output CycB $= \neg$cdh1 $\wedge$ $\neg$Cdc20---only
  two effective regulators.
\item G9: Output CHKREC has 8 regulators, but a null-state clause
  (CHKREC $= 1$ when all damage genes are off) creates a large
  degenerate region in the value function.
\end{itemize}
%
G4 (Davidich fission yeast) is nominally negative
($\Delta_{3+} = -0.4$~pp) but its effect is $8\times$ smaller than
any other network's. The output is gated by Ste9, Rum1, and Slp1,
but Ste9 and Rum1 have identical update rules, reducing effective
gating to two variables. This functional redundancy may explain both
the weak local higher-order structure and the near-zero composition
gap.


\subsection{Exploratory: limit-cycle prevalence and gap direction}
\label{sec:cycling}

Rank-ordering the six networks by cycling fraction places all three
creation networks ($\Delta_{3+} > 0$) below all networks with
$\Delta_{3+} \leq 0$ (Table~\ref{tab:results}, Cycling column).
The Spearman correlation between $\Delta_{3+}$ and cycling fraction
is $\rho = -0.83$ ($p = 0.04$). Under a permutation test on the sign
pattern alone, the probability that three positives precede three
non-positives among six ordered items is $1/\binom{6}{3} = 0.05$
one-tailed, but because this predictor was identified after the data
were collected (the pre-registered predictor, fb\_ratio, failed;
Section~\ref{sec:blind-results}), the relevant $p$-value is
two-tailed: $p \approx 0.10$.

Two further caveats weaken the pattern. First, the boundary case is
G4 ($\Delta_{3+} = -0.4$~pp, cycling $= 17.7\%$), whose effect
is effectively zero (Section~\ref{sec:results}). If G4 is treated as
a null rather than a destruction event, the sign separation becomes
3/2/1 and the threshold has no clean cut point. Second, magnitude is
not monotone in cycling: G10 (cycling $= 16.4\%$) shows $4\times$ the
creation of G2 (cycling $= 7.4\%$), so there is no dose-response
relationship.

A plausible mechanism exists: time-averaging over limit-cycle periods
smooths coalition-dependent variation in attractor structure,
compressing higher-order Walsh coefficients. Fixed-point attractors
preserve discrete switching behavior, allowing combinatorial
interactions to compound. Cycling fraction is weakly correlated with
network size ($\rho = 0.09$, $p = 0.87$) and moderately with mean
in-degree ($\rho = 0.54$, $p = 0.27$), so cycling is not a simple
proxy for either.

This association is a hypothesis for future testing on additional
networks, not a confirmed result.


\subsection{Pre-registered structural predictions}
\label{sec:blind-results}

Table~\ref{tab:blind} summarizes the blind prediction scorecard.

\begin{table}[t]
\centering
\caption{Blind prediction accuracy. Predictions were registered
before any coalition sweep ran, based solely on structural analysis
of Boolean rules and feedback topology.}
\label{tab:blind}
\begin{tabular}{lcc}
\toprule
Category & Score & Notes \\
\midrule
$\rho$ direction (positive) & 6/6 & All predicted positive, all observed positive \\
$\rho$ in predicted range & 4/6 & G4 below range; G9 slightly above \\
Gap direction & 4/6 & Missed G1, G4 (overweighted positive feedback) \\
Spectrum in range & 3/6 & G1, G9 below; G10 far above \\
\bottomrule
\end{tabular}
\end{table}

\textbf{What structural analysis gets right.}
The universal positivity of $\rho$ was correctly predicted from first
principles: local interactions must contribute \emph{something} to
global interactions because each gene's update rule is one component
of the composed dynamical system. The gap direction was correctly
predicted in 4 of 6 networks, including the only destruction
prediction (G9, based on negative feedback dominance and the CHKREC
global reset mechanism).

\textbf{What structural analysis gets wrong.}
The two misses (G1, G4) both involved overweighting positive feedback
loop counts. The agent predicted that strong positive feedback creates
bistability and higher-order interactions. In G4, the strongest
positive feedback ratio in the panel (13:3) instead creates a
binary-switch attractor that compresses the value function. Positive
feedback can linearize as easily as it can complexify.

\textbf{Cross-network hypothesis: falsified.}
The pre-registered hypothesis---fb\_ratio $> 0.5$ implies
creation---fails decisively. G4 (fb\_ratio $= 0.81$) shows
destruction; G10 (fb\_ratio $= 0.48$) shows the strongest creation.
Feedback loop counting does not predict the composition gap.


%% ====================================================================
\section{Discussion}
\label{sec:discussion}
%% ====================================================================

\subsection{Implications for knockout screen design}

The molecular wiring diagram is a statistically significant predictor
of pairwise knockout epistasis in three of six networks ($\rho =
0.25$--$0.41$, $p < 0.05$). In the other three, local rule structure
provides no useful prior for selecting which gene pairs to test. This
heterogeneity means that the common practice of using pathway
databases to prioritize combinatorial experiments is justified in some
regulatory contexts but misleading in others.

The composition gap provides a more nuanced guide. In creation
networks (G2, G5, G10), dynamics generate interactions absent from
any local rule. A knockout screen guided by molecular interactions
alone would miss the strongest epistatic pairs. In destruction
networks (G1, G9), the molecular wiring diagram over-predicts
higher-order interactions that dynamics compress. A screen guided by
local rules would over-allocate resources to combinations whose
epistasis is damped by dynamical convergence.

\subsection{Relationship to prior work}

Gjuvsland et al.\ \citeneeded{Gjuvsland 2007} demonstrated that
statistical epistasis is generic in gene regulatory networks. Our
results extend this in two directions. First, the gap is quantified
order-by-order rather than detected as present/absent. Second, the
gap is signed: creation and destruction are distinct phenomena with
different structural correlates. A pre-registered attempt to predict
the gap direction from feedback loop topology failed, which is itself
informative: it demonstrates that the composition gap is not
reducible to simple topological features of the regulatory graph.

Baier et al.\ \citeneeded{Baier 2023} measured epistasis in a
synthetic three-gene network, finding environment-dependent sign
switching. Our Walsh decomposition provides a natural framework for
their observation: the composition gap at each order can change sign
across conditions (clamp-to-0 vs.\ clamp-to-1), and we anticipate
that extending our analysis to both ablation baselines will reveal
signed sensitivity at the order level.

The global epistasis literature \citeneeded{Otwinowski 2018, Reddy
Desai 2021} establishes that nonlinear genotype-phenotype maps
reshape interaction spectra. Our local-vs-global comparison
operationalizes this reshaping as a quantitative spectral
transformation, decomposed by interaction order.


\subsection{Limitations}

\begin{itemize}
\item $N = 6$ networks. All cross-network associations (cycling
  fraction, output complexity) are descriptive case studies.
  Confirmatory regression requires $N \geq 15$--$20$, and
  \emph{Arabidopsis} alone drives most of the variance in creation
  magnitude.
\item Clamp-to-0 baseline only. Clamp-to-1 may yield different gap
  directions, as the asymmetry of Boolean logic (AND vs.\ OR) makes
  the value function baseline-dependent.
\item All genes are players. Excluding constitutive input nodes
  (which have no regulators) from the player set would change the
  pairwise space and potentially sharpen correlations.
\item Synchronous update only. Asynchronous dynamics would change
  attractor structure and cycling fractions.
\item Sum-aggregation for local pairwise magnitudes
  (Eq.~\ref{eq:local-pair}). Max-aggregation is an untested
  alternative.
\item $n \leq 15$. The exhaustive approach does not scale to
  genome-wide networks. These results characterize the composition
  gap in well-studied regulatory modules, not whole-genome epistasis.
\end{itemize}


%% ====================================================================
\section{Conclusion}
\label{sec:conclusion}
%% ====================================================================

The gap between molecular interaction structure and phenotypic
epistasis is quantifiable, signed, and large. In six Boolean gene
regulatory networks, we measured the complete Walsh interaction
spectrum at both the local rule and attractor levels. The composition
gap spans two orders of magnitude, from effectively zero to 34
percentage points, and its sign varies across network architectures.
The molecular wiring diagram provides a universally positive but
often non-significant prior for predicting knockout epistasis. A
pre-registered hypothesis based on feedback loop counting was
falsified; an exploratory association with limit-cycle prevalence
awaits validation on additional networks.

The complete Walsh spectra, coalition value functions, and
pre-registration artifacts are available at
\textcolor{red}{[URL: repository]}.


%% ====================================================================
%% References
%% ====================================================================

\bibliographystyle{plainnat}
% \bibliography{references}

\end{document}
